\documentclass[11pt]{article}
\usepackage[margin=1in]{geometry}
\usepackage{amsmath,amssymb,amsthm,mathtools,booktabs}
\usepackage[T1]{fontenc}
\usepackage{lmodern,microtype}
\usepackage{xcolor}
\usepackage[colorlinks=true,linkcolor=blue!45!black,urlcolor=blue!45!black,citecolor=blue!45!black]{hyperref}
\newtheorem{theorem}{Theorem}[section]
\newtheorem{lemma}[theorem]{Lemma}
\newtheorem{corollary}[theorem]{Corollary}
\newtheorem{proposition}[theorem]{Proposition}
\theoremstyle{remark}\newtheorem{remark}[theorem]{Remark}
\newcommand{\R}{\mathbb R}\newcommand{\C}{\mathbb C}\newcommand{\dd}{\,\mathrm d}
\newcommand{\ee}{\mathrm e}\newcommand{\abs}[1]{\lvert#1\rvert}
\newcommand{\GRH}{\mathrm{DGRH}}
\newcommand{\spec}{\mathcal F}
\newcommand{\FT}{\widehat\mu}
\DeclareMathOperator{\ord}{ord}
\title{Positive filters for the full Dirichlet zero spectrum\\\large A sharp double-exponential localization theorem and exact prime-error averaging}
\author{Research note for ProveIt}
\date{27 November 2026}
\hypersetup{pdftitle={Positive filters for the full Dirichlet zero spectrum},pdfauthor={Research note for ProveIt}}
\setlength{\emergencystretch}{2em}
\begin{document}
\maketitle
\begin{abstract}
The merged ProveIt article constructs compactly supported positive filters cancelling finitely many Dirichlet zero ordinates. We pass to a single positive filter cancelling every noncentral ordinate, with aggregate multiplicity. Assuming Dirichlet GRH, for a conjugation-closed multiset of $d$ primitive Dirichlet $L$-functions (including the zeta function), there exists an even, strictly positive probability density $p$ with
$p(t)\ll \exp(-c\exp(2|t|/d))$ whose Fourier transform has all these prescribed zeros. Conversely no probability measure with this zero divisor can have a finite moment $\int\exp(c\exp(\alpha|t|))\,\dd\mu(t)$ when $\alpha>2/d$. The exponent $2/d$ is thus exact. The construction uses central-zero division of completed $L$-functions, an eventual-positivity argument for convolution of theta kernels, and a general compact-smoothing lemma for signed densities. As an arithmetic consequence, one positive rule cancels the complete oscillatory spectrum of all the associated smooth prime errors, leaving the central-zero term and an $O(1)$ remainder in the original scale. Truncating the rule retains this remainder with logarithmic half-width $(d/2)\log\log x+O(1)$. No independence or simplicity of zeros is required.
\end{abstract}

\section{Input, scope, and comparison with the article}
The input is the Dirichlet GRH exported by \texttt{grh-lean}, vendored in ProveIt at upstream revision
\begin{center}\small\texttt{72fead102f663d5c7a407964a13dc0cf9d977207}.\end{center}
We inspected the two repository directories specified in the request, the actual statements and main proof interface, and the relevant portions of the merged article \cite{vendor,article}. The statement file imports only Mathlib. The precise half-plane theorem in \texttt{src/Main.lean} is
\begin{quote}\small\texttt{LFunction\_ne\_zero\_of\_half\_lt\_re : DirichletLNonvanishingRightOfHalf},\end{quote}
where the type quantifies over every positive modulus, every Dirichlet character and every complex argument with real part greater than $1/2$, except the principal pole. The theorem
\texttt{LFunction\_eq\_zero\_re} puts every zero in the open critical strip on its middle line. The file \texttt{src/New/Main.lean} uses the functional equation to pass from the right-half-plane theorem to the critical-line theorem. The file \texttt{src/New/Reduction/OneSidedBound.lean} connects a one-sided smoothed prime remainder to nonvanishing. Its opening lemma bounds the remainder at one scale by an arithmetic main-inequality contribution and remainders at a much larger scale. The package reports successful independent kernel checks and only the usual Lean axioms; this note has not rebuilt or replayed the proof. \texttt{Challenge.lean} is a separate comparator specification, not the proof entry point. The source package pins Lean 4.35.0 and its own Mathlib revision; the ProveIt root build is not a check of that package.

The merged article's \texttt{rec:thm:finitefilter} and \texttt{rec:thm:logfilter} cancel a finite multiset of frequencies with a positive compactly supported kernel of half-width $O_{\spec}(\log B)$ for zero height $B$. Its \texttt{rec:lem:jensensupport} gives a lower bound of order $\log B$. Its \texttt{rec:thm:uniformavg} and \texttt{rec:thm:polylogavg} use those finite filters for uniform approximate arithmetic averaging. Our question instead requires exact simultaneous cancellation at \emph{every} noncentral zero, so compact support is impossible. We determine the sharp double-exponential tail exponent for a positive filter on the whole line, including complex characters and central zeros. The article's source description and the relevant theorem statements are directly read from LaTeX source, not inferred from its PDF or title. See \texttt{SOURCE\_AUDIT.md} for the scope and blob identifiers. The article describes itself as a merge rather than an independent verification.

The classical Riemann $\Xi$ Fourier kernel is already positive, and is a known characteristic-function example \cite{titchmarsh,polson}. We do not claim that example, theta inversion, Jensen's formula or the explicit formula as new. The proposed contribution is the general signed-to-positive compact smoothing construction, its use for arbitrary conjugation-closed Dirichlet families after removing their central zeros, the matching tail obstruction, and the arithmetic consequence. Targeted comparison with the article and literature does not establish priority throughout the mathematical literature; the novelty claim is provisional. All statements below are proved mathematically, but not machine-formalized here.

\section{The full-spectrum theorem}
Let $\spec$ be a nonempty finite multiset of primitive characters closed under complex conjugation; the trivial character of conductor $1$ is allowed. Repetitions count, and $d=|\spec|$. For $\chi\ne\mathbf1$ let $q_\chi$ be its conductor and $a_\chi\in\{0,1\}$ its parity. Write
\begin{align*}
\Lambda_\chi(s)&=(q_\chi/\pi)^{(s+a_\chi)/2}\Gamma((s+a_\chi)/2)L(s,\chi),\\
\Lambda_{\mathbf1}(s)&=\xi(s)=\tfrac12 s(s-1)\pi^{-s/2}\Gamma(s/2)\zeta(s).
\end{align*}
We use the classical analytic continuation, functional equations, and zero counting for these completions \cite{davenport,titchmarsh}. Their zeros are the nontrivial zeros; the zeta poles have been removed. Let $m_\chi(\gamma)$ be the multiplicity of $1/2+i\gamma$ in the completion and $m(\gamma)=\sum_{\chi\in\spec}m_\chi(\gamma)$. Under $\GRH$ these exhaust the zeros. Our Fourier convention is $\widehat\mu(z)=\int\ee^{izt}\dd\mu(t)$.

\begin{theorem}[Sharp full-spectrum localization]\label{thm:main}
Assume $\GRH$. There is an even, strictly positive continuous probability density $p$ on $\R$ and constants $C,c>0$ such that
\begin{equation}\label{eq:tail}
 p(t)\le C\exp\big(-c\exp(2|t|/d)\big),\qquad t\in\R,
\end{equation}
and $\widehat p$ is entire with $\ord_\gamma\widehat p\ge m(\gamma)$ for every $\gamma\ne0$ with $m(\gamma)>0$. Conversely, if any probability measure $\mu$ has these zeros with at least the prescribed multiplicities and
\begin{equation}\label{eq:moment}
 \int_\R \exp\big(b\exp(\alpha|t|)\big)\dd\mu(t)<\infty
\end{equation}
for some $b,\alpha>0$, then $\alpha\le 2/d$. Thus the maximum achievable moment exponent is exactly $2/d$ (including its endpoint). Additional zeros of $\widehat p$ are allowed. Central zeros are deliberately excluded because $\widehat p(0)=1$.
\end{theorem}

An explicit construction, expressed using convergent integrals and a compact smoothing kernel, is given below. The constants depend on the family; no bound uniform in conductor is asserted.

\subsection{A compact-smoothing lemma}
\begin{lemma}[Signed-to-positive compact smoothing]\label{lem:smooth}
Let $h:\R\to\R$ be continuous and even, strictly positive for $|t|>A$ for some $A$, and satisfy $\int |h(t)|\ee^{u|t|}\dd t<\infty$ for every $u\ge0$. Suppose its bilateral Laplace transform $H(u)=\int h(t)\ee^{ut}\dd t$ is strictly positive for every real $u$. Then there exist $\sigma,R>0$ such that
\[
 k(t)=\exp(-t^2/(2\sigma^2))\mathbf1_{[-R,R]}(t),\qquad h*k>0\quad\text{on }\R.
\]
In particular this convolution can be normalized to an even positive probability density. If $|h(t)|\le C\exp(-c\exp(\alpha |t|))$, the normalized convolution obeys a bound of the same form with exponent $\alpha$. Conversely, convolution with a nonzero nonnegative compactly supported integrable kernel cannot produce an integrable nonzero nonnegative function if $H$ vanishes anywhere on the real axis.
\end{lemma}
\begin{proof}
Enlarge $A$ so $A\ge1$. Put $N=\int h_-$ and $P=\int_{A+1}^{A+2}h(t)\dd t>0$, and choose $D>0$ with $\ee^D P>N$. For $|u|\le D+1$ the continuous function $H(u)$ has positive minimum $\delta$. Choose $B>A+2$ so that
\[
 \int_{|t|>B}|h(t)|\ee^{(D+1)|t|}\dd t<\delta/4.
\]
Uniform convergence on the compact interval of $u$ shows that for sufficiently large $\sigma$, with $\sigma^2\ge A+2$,
\begin{equation}\label{eq:truncpos}
 \int_{-B}^B h(t)\ee^{ut}\ee^{-t^2/(2\sigma^2)}\dd t>\delta/2\quad (|u|\le D+1).
\end{equation}
Take $R=D\sigma^2+A+2+B$. For $|x|\le D\sigma^2+A+2$, the interval $|t|\le B$ lies wholly in the region $|x-t|\le R$, and $|x/\sigma^2|\le D+1$. Factor $\ee^{-x^2/(2\sigma^2)}$ out of the convolution on this interval and apply \eqref{eq:truncpos} with $u=x/\sigma^2$. All omitted contributions are nonnegative because $B>A$. Hence $(h*k)(x)>0$ throughout the middle region.

For $x\ge D\sigma^2+A+2$, every negative contribution has $t\le A$ and hence $k(x-t)\le k(x-A)$, by monotonicity in the nonnegative argument. If $k(x-A)>0$, then for $t\in[A+1,A+2]$ we have $0\le x-t\le x-A\le R$ and
\[
 k(x-t)\ge k(x-A-1)=k(x-A)\exp((x-A-1/2)/\sigma^2)\ge\ee^D k(x-A).
\]
Thus the convolution is at least $k(x-A)(\ee^DP-N)>0$. If $k(x-A)=0$, all negative contributions vanish, while a neighborhood of $t=x$ gives a strictly positive contribution. Evenness handles $x$ negative. The convolution is continuous by standard $L^1$ convolution/translation continuity (here $h$ is continuous and locally bounded and the other factor has compact support). The integral is $H(0)\int k>0$. For the tail, $|t-v|\ge |t|-R$ for $|v|\le R$, so one can replace $c$ by $c\ee^{-\alpha R}$ and multiply $C$ by a constant. Finally, the Laplace transform of a convolution is the product; the Laplace transform of a nonzero nonnegative function is strictly positive wherever it is finite. This proves necessity.
\end{proof}

\subsection{Theta kernels and their convolutions}
Define $F(w)=\prod_{\chi\in\spec}\Lambda_\chi(1/2+w)$. We prove that the inverse bilateral Laplace kernel of $F$ is real, eventually strictly positive on the right, and has decay exponent $2/d$. For a nontrivial primitive character use
\begin{equation}\label{eq:theta}
 K_\chi(t)=2\ee^{(1/2+a_\chi)t}\sum_{n\ge1}\chi(n)n^{a_\chi}\exp(-\pi n^2\ee^{2t}/q_\chi).
\end{equation}
For $\Re(s)>1$ the Mellin integral, termwise justified by absolute convergence, gives $\int K_\chi(t)\ee^{wt}\dd t=\Lambda_\chi(1/2+w)$. Theta inversion for primitive characters gives $K_\chi(-t)=\varepsilon_\chi K_{\bar\chi}(t)$, with $|\varepsilon_\chi|=1$ (equivalently use the usual functional equation). This implies decay at the other end; the integral is entire and hence agrees everywhere by continuation. On the right the $n=1$ term is real and positive and all the others are relatively $O_\chi(\exp(-c_\chi\ee^{2t}))$. Indeed, after factoring out that term, for large $t$ the sum of $n^{a_\chi}\exp(-\pi(n^2-1)\ee^{2t}/q_\chi)$ over $n\ge2$ is dominated by a convergent Gaussian series times $\exp(-c_\chi\ee^{2t})$.

For conductor $1$ let $g(t)=2\ee^{t/2}\sum_{n\ge1}\exp(-\pi n^2\ee^{2t})$ and set $K_{\mathbf1}=(g''-g/4)/2$. Explicitly,
\begin{equation}\label{eq:xi}
 K_{\mathbf1}(t)=\ee^{t/2}\sum_{n\ge1}(4A_n^2-6A_n)\ee^{-A_n},\qquad A_n=\pi n^2\ee^{2t}.
\end{equation}
The classical theta identity gives $g(-t)=g(t)+\ee^{t/2}-\ee^{-t/2}$, so $K_{\mathbf1}$ is even. Integrating by parts for $\Re w>1/2$ and then continuing shows its transform is $\xi(1/2+w)$, since $(w^2-1/4)/2=s(s-1)/2$. The boundary terms vanish in that half-plane by the theta identity and rapid right decay. All summands of \eqref{eq:xi} are positive for $t\ge0$, and the first dominates there as $t\to\infty$. These are standard formulas; the rest of the argument uses their tail properties rather than positivity of any individual nontrivial-character kernel.

In particular for all $d$ factors, indexed with repetition by $j$, there are common constants $C,c>0$ with
\begin{equation}\label{eq:Kbound}
 |K_j(t)|\le C\exp(-c\exp(2|t|)),\qquad
 K_j(t)=v_j(t)(1+o(1))\quad(t\to+\infty),
\end{equation}
where $v_j(t)>0$ and $v_j(t)\ge\exp(-C\exp(2t))$ for sufficiently large $t$. The relative error may be complex. The linear exponential prefactors and the zeta polynomial prefactor are absorbed into the constants. This uses only classical theta identities, not GRH.

\begin{lemma}[Convolution tails]\label{lem:tails}
The function $f=K_1*\cdots*K_d$ is continuous, real, has bilateral Laplace transform $F$, and satisfies
\[
 |f(t)|\le C_1\exp(-c_1\exp(2|t|/d)),\qquad f(t)>0\quad(t\text{ sufficiently large}).
\]
There exists a sign $\eta\in\{-1,1\}$ with $f(-t)=\eta f(t)$.
\end{lemma}
\begin{proof}
The transform identity follows from absolute convergence and Fubini. Conjugation of the family makes $f$ real (the kernels of conjugate characters are conjugate), and the product functional equation gives $F(-w)=\eta F(w)$, where the product of the root numbers is real and of modulus one. Fourier uniqueness yields the parity. Continuity follows from the continuous integrable bounded kernels. For the tail estimate write $t_1+\cdots+t_d=t$ and integrate over the first $d-1$ coordinates, with $t_d$ determined. By convexity,
\[
 \sum_{j=1}^d\ee^{2|t_j|}\ge d\ee^{2|t|/d}.
\]
Split the exponent in \eqref{eq:Kbound} into two equal parts; the first produces the claimed tail bound, and the remaining integral is bounded uniformly by
$\big(\int_\R\exp(-c\ee^{2|u|}/2)\dd u\big)^{d-1}$ (discard the last factor). The case $d=1$ follows directly from \eqref{eq:Kbound}.

To prove eventual positivity for $d\ge2$, choose $T$ large enough that whenever all $t_j\ge T$, the real part of the product of the kernels is at least $\frac12\prod_j v_j(t_j)$, using the relative errors in \eqref{eq:Kbound} and fixed $d$. The integral over this region has nonnegative real part. On the box $|t_j-t/d|\le 1/d$ for $j<d$, the last coordinate also lies within $1$ of $t/d$. For large positive $t$, this gives a contribution at least $c_2\exp(-C_2\ee^{2t/d})$, with $c_2>0$. On the complement some $t_j<T$, so
\[
 \sum_{i\ne j}\ee^{2|t_i|}\ge(d-1)\ee^{2(t-T)/(d-1)}.
\]
Again split the exponent in two and integrate the remaining factor as before. The total absolute contribution from this complement is $O_d(\exp(-c_3\ee^{2(t-T)/(d-1)}))$, negligible compared with the lower bound. Since $f$ is real, it is eventually strictly positive. For $d=1$ the family is self-conjugate, and the right-tail asymptotic already proves the assertion.
\end{proof}

\subsection{Removing the central zero and proving existence}
Let $m=\ord_{w=0}F(w)$, a finite nonnegative integer. Define the entire function $H(w)=F(w)/w^m$ with the removable value at zero. If $m=0$, let $h=f$; otherwise put
\begin{equation}\label{eq:antiderivative}
 h(t)=\frac1{(m-1)!}\int_t^\infty (v-t)^{m-1}f(v)\dd v.
\end{equation}
Because $F^{(j)}(0)=\int v^j f(v)\dd v=0$ for $0\le j<m$, the same integral may be written as minus the integral from $-\infty$ to $t$ with the same integrand. This identity gives decay on the left, while \eqref{eq:antiderivative} gives decay on the right. For example for $t\ge0$ split the decay exponent in Lemma~\ref{lem:tails} into two halves, pull out $\exp(-c_1\ee^{2t/d}/2)$, and integrate the remaining bound against $v^{m-1}$ after substituting $v-t$ for $v$. The left tail follows in exactly the same way using the alternative formula and $t-v\ge0$. Thus $h$ has a tail bound of the form \eqref{eq:tail}. Its first $m-1$ derivatives have the analogous bounds, and $h^{(m)}=(-1)^m f$. Integration by parts $m$ times, with vanishing boundary terms for every $w$, now gives $\int h(t)\ee^{wt}\dd t=H(w)$. The function $h$ is real and continuous. As $\eta=(-1)^m$ by parity of the first nonzero Taylor coefficient of $F$, $H$ is even, and Fourier uniqueness implies $h$ is even. Finally $h$ is positive sufficiently far to the right by the eventual positivity of $f$ and the nonnegative integrand multiplier in \eqref{eq:antiderivative}; thus it is positive also sufficiently far to the left.

Under Dirichlet GRH, $F(w)$ has no real zero except possibly at $w=0$, and hence $H(u)$ is nowhere zero on the real axis. The pole-removed completions are nonzero at $s=0,1$ by the standard functional equations and the nonvanishing of $L(1,\chi)$; these points introduce no extra exception. For positive $u$ sufficiently large, $F(u)>0$: the factors for complex conjugates multiply to positive absolute squares, the real-character Euler products are positive, and the zeta completion is positive there. Consequently $H(u)>0$ on the whole real axis. Apply Lemma~\ref{lem:smooth} to obtain $k$ and set
\begin{equation}\label{eq:construction}
 p(t)=\frac{(h*k)(t)}{H(0)\int_\R k(v)\dd v}.
\end{equation}
This is a strictly positive even probability density with the required tail bound. Its Fourier transform is
\[
 \widehat p(z)=\frac{H(iz)}{H(0)}\cdot\frac{\int k(v)\ee^{izv}\dd v}{\int k(v)\dd v}.
\]
At nonzero $z=\gamma$ the central-zero division is harmless, and $F(i\gamma)$ vanishes to the required aggregate order. Multiplication by an entire transform cannot remove zeros. The existence part of Theorem~\ref{thm:main} follows. For any $0<b<c$ the double-exponential moment in \eqref{eq:moment} at $\alpha=2/d$ is finite by integrating \eqref{eq:tail}.

\begin{remark}[What part of GRH is needed for positivity?]
The construction through $h$ is unconditional. In fact the existence of a nonnegative compact smoothing of this particular central-zero-removed kernel is equivalent to $F$ having no real noncentral zeros. Necessity is the last assertion of Lemma~\ref{lem:smooth}; sufficiency is the preceding argument using the sign at large positive $u$. Full GRH is needed to identify the entire nontrivial zero divisor with the real Fourier frequencies in our prime-counting application. This equivalence does not assert a new RH for other classes of functions.
\end{remark}

\subsection{The lower bound: a sharp Jensen obstruction}
Suppose a probability measure $\mu$ obeys \eqref{eq:moment}, with value $M<\infty$. Its Fourier transform is entire by domination on every compact set, and $\widehat\mu(0)=1$. For $z=R\ee^{i\theta}$, set $y=R|\sin\theta|$. Then
\[
 |\widehat\mu(z)|\le M\exp\left(\sup_{r\ge0}\{yr-b\ee^{\alpha r}\}\right).
\]
When $y\ge b\alpha$, the supremum equals $(y/\alpha)(\log(y/(b\alpha))-1)$, and when $y<b\alpha$ it is $-b$. In particular, uniformly in $\theta$ for $R\ge2$,
\begin{equation}\label{eq:growth}
 \log|\widehat\mu(R\ee^{i\theta})|\le \frac{R|\sin\theta|}{\alpha}\log R+O_{b,\alpha,M}(R).
\end{equation}
Jensen's formula, at radii avoiding boundary zeros and then by limits, gives
\begin{equation}\label{eq:jensen}
 \sum_{0<|z|<R\, :\,\widehat\mu(z)=0}\ord_z\widehat\mu\,\log(R/|z|)
 \le \frac{2}{\pi\alpha}R\log R+O(R).
\end{equation}
For the fixed family the classical Riemann--von Mangoldt formula, summed with multiplicity, gives
\[
 N_+(T)=\sum_{0<\gamma\le T}m(\gamma)=\frac{d}{2\pi}T\log T+O_{\spec}(T).
\]
For complex characters one may count positive and negative ordinates together and use conjugation symmetry of the family to halve the count. Under GRH these are exactly the relevant zeros. The finitely many central zeros do not contribute. Keeping the prescribed positive and negative real zeros in \eqref{eq:jensen} yields
\[
 2\int_0^R\frac{N_+(u)}u\dd u=\frac{d}{\pi}R\log R+O_{\spec}(R).
\]
The integral at zero is finite because the nonzero zeros of a nonzero entire function are isolated. Comparing leading coefficients proves $d\le 2/\alpha$, as claimed. No assumption on distinctness, simplicity, gaps, or independence of ordinates occurs.

\begin{corollary}[Endpoint constant obstruction]\label{cor:constant}
Write $Q_*=(\prod_{\chi\in\spec}q_\chi)^{1/d}$, with $q_{\mathbf1}=1$. If \eqref{eq:moment} holds with $\alpha=2/d$, then necessarily $b\le 2\pi d/Q_*$. This is a necessary bound, not a claim of sharpness for the constant $b$.
\end{corollary}
\begin{proof}
Use the refined zero count
\[
 N_+(T)=\frac{dT}{2\pi}\left(\log T+\log(Q_*/(2\pi))-1\right)+O_{\spec}(\log(T+2)).
\]
Integration gives the prescribed Jensen sum $\frac{dR}{\pi}(\log R+\log(Q_*/(2\pi))-2)+O_{\spec}(\log^2(R+2))$. On the upper side retain the exact supremum above: its small-$y$ branch changes the angular average by $O(1)$ (extend $y\log y$ continuously at zero). The mean of $|\sin\theta|$ is $2/\pi$, and that of $|\sin\theta|\log|\sin\theta|$ is $(2/\pi)(\log2-1)$, as follows by substituting $u=\cos\theta$ on a quadrant. Hence the upper Jensen mean is
\[
 \frac{2R}{\pi\alpha}\left(\log R-\log(b\alpha)+\log2-2\right)+O(1).
\]
At $\alpha=2/d$ compare the coefficients of $R$ to obtain $\log(Q_*/(2\pi))\le \log2-\log(2b/d)$, equivalently $b\le2\pi d/Q_*$.
\end{proof}

\section{Arithmetic consequence: one exact spectral averaging rule}
\label{sec:arithmetic}
Let $W\in C_c^\infty((1,2))$ be real, with Mellin transform $\widetilde W(s)=\int_0^\infty W(v)v^{s-1}\dd v$. For a primitive character define the smoothed sum
\[
 S_\chi(x)=\sum_{n\ge1}\Lambda(n)\chi(n)W(n/x),\qquad x>0,
\]
where the unadorned $\Lambda(n)$ is the von Mangoldt function. Let $\delta_\chi=1$ for conductor $1$ and $0$ otherwise, and define the normalized error
\[
 E_\chi(t)=\ee^{-t/2}\big(S_\chi(\ee^t)-\delta_\chi\ee^t\widetilde W(1)\big).
\]
The smooth explicit formula \cite{davenport,titchmarsh} and GRH give, for $t\ge0$,
\begin{equation}\label{eq:explicit}
 E_\chi(t)=-\sum_{\gamma}m_\chi(\gamma)\widetilde W(1/2+i\gamma)\ee^{i\gamma t}+r_\chi(t),\qquad
 |r_\chi(t)|\ll_{\chi,W}\ee^{-t/2}.
\end{equation}
Here the series includes $\gamma=0$ and is absolutely and uniformly convergent: repeated integration by parts in $\log v$ gives rapid decay of $\widetilde W(1/2+i\gamma)$, and the zero count grows as $O(T\log(T+2))$. For completeness, Mellin inversion on a line $\Re s>1$ expresses $S_\chi$ as the integral of $-L'/L(s,\chi)x^s\widetilde W(s)$. Shifting to the left past the pole and zeros gives the displayed nontrivial-zero series plus the trivial-zero terms. In the smooth explicit formula the trivial zeros lie at $-a_\chi-2j$, $j\ge0$ for nonprincipal primitive characters, and at $-2j$, $j\ge1$ for zeta. Their total is $O_{\chi,W}(1)$ for $x\ge1$, because the support stays a positive distance from $1$ and $|\widetilde W(-r)|\ll_W (1+\epsilon_W)^{-r}$ for some $\epsilon_W>0$. Dividing by $\sqrt x$ yields the bound on $r_\chi$. The standard contour shift uses horizontal heights away from zeros, on which the logarithmic derivative has polynomial bounds, and the rapidly decaying smooth Mellin factor; the left vertical contours can be taken between consecutive trivial zeros. For $t<0$ define $r_\chi$ by the same identity. The sum defining $S_\chi(\ee^t)$ vanishes for $t$ sufficiently negative, and $E_\chi$ is bounded on $(-\infty,0]$. The spectral series is uniformly bounded as well, so after enlarging the constant,
\begin{equation}\label{eq:globalr}
 |r_\chi(t)|\le C_{\chi,W}\ee^{-t/2}\quad(t\in\R),\qquad \sup_{t\in\R}|E_\chi(t)|<\infty.
\end{equation}

\begin{theorem}[Full cancellation with bounded unnormalized error]\label{thm:prime}
Take $p$ from Theorem~\ref{thm:main}, let $\mu$ be the corresponding probability measure, and define $M_\mu(v)=\int\ee^{vu}\dd\mu(u)$. For every $\chi\in\spec$ and every such $W$,
\begin{align}
 \int E_\chi(t+u)\dd\mu(u)&=-m_\chi(0)\widetilde W(1/2)+O_{\spec,W}(\ee^{-t/2}),\label{eq:normalized}\\
 \int\ee^{-u/2} S_\chi(x\ee^u)\dd\mu(u)&=\delta_\chi x\widetilde W(1)M_\mu(1/2)
 -\sqrt x\,m_\chi(0)\widetilde W(1/2)+O_{\spec,W}(1).\label{eq:unnorm}
\end{align}
The rule depends only on the family, not on $W$ or $x$. The error bounds hold uniformly as $t\to\infty$ or $x\to\infty$. The zero-dependent oscillatory part in \eqref{eq:explicit} is cancelled exactly, not merely up to a height cutoff.
\end{theorem}
\begin{proof}
By absolute summability and Fubini the term at ordinate $\gamma$ in the average in \eqref{eq:normalized} is multiplied by $\widehat\mu(\gamma)$, zero for all nonzero ordinates and one at zero. The remainder is at most $C_{\chi,W}\ee^{-t/2}M_\mu(-1/2)$ by \eqref{eq:globalr}; this moment is finite. Multiply \eqref{eq:normalized} by $\sqrt x$ with $x=\ee^t$ and insert the main term in $E_\chi$ to obtain \eqref{eq:unnorm}. The family is finite, so the constants can be made simultaneous. All integrals converge absolutely by \eqref{eq:globalr} and the exponential moments of $\mu$.
\end{proof}

\begin{corollary}[A finite polylogarithmic dilation range with explicit exponent]\label{cor:trunc}
For sufficiently large $x$, set
\[
 H_x=\frac d2\log\log x+C_0,
 \qquad \dd\mu_x(u)=\frac{\mathbf1_{|u|\le H_x}\dd\mu(u)}{\mu([-H_x,H_x])},
\]
where $C_0$ is a sufficiently large family-dependent constant. Formula \eqref{eq:unnorm} remains true with $\mu_x$ on the left and even with the original $M_\mu(1/2)$ on the right. Thus only the dilations
\[
 \ee^{-C_0}x(\log x)^{-d/2}\le x\ee^u\le \ee^{C_0}x(\log x)^{d/2}
\]
are needed. This is an upper bound on the width of a truncated approximate filter; no matching lower bound for this \emph{approximate} task is asserted.
\end{corollary}
\begin{proof}
From \eqref{eq:tail}, for each fixed real $v$,
\[
 \int_{|u|>H}\ee^{v u}\dd\mu(u)\ll_{\spec,v}\exp(-c'\exp(2H/d))
\]
for some $c'>0$. One obtains this, for example, by allocating half of the double-exponential factor to dominate $\ee^{|v||u|}$ integrably. Choose $C_0$ so that the bound is $O(x^{-2})$ for $v=0,1/2,-1/2$. Renormalization changes averages of the bounded $E_\chi$ by $O(x^{-2})$, and changes $M_\mu(1/2)$ by $O(x^{-2})$. The changes in \eqref{eq:unnorm} are $O(x^{-3/2})$ and $O(x^{-1})$, respectively. The original $O(1)$ remainder is unchanged.
\end{proof}

\begin{corollary}[All reduced residue classes at once]\label{cor:residues}
Fix a modulus $q\ge1$. Let $\spec$ consist of the primitive characters inducing all the characters modulo $q$; then $d=\varphi(q)$. For $a\in(\mathbb Z/q\mathbb Z)^\times$ put
\[
 S_a(x)=\sum_{n\equiv a\bmod q}\Lambda(n)W(n/x),\qquad
 B_a=\frac1{\varphi(q)}\sum_{\chi\bmod q}\overline{\chi(a)}\,m_{\chi^*}(0),
\]
where $\chi^*$ denotes the inducing primitive character. The same positive $\mu$ works for every $a$ and $W$, and
\[
 \int\ee^{-u/2}S_a(x\ee^u)\dd\mu(u)
 =\frac{x\widetilde W(1)}{\varphi(q)}M_\mu(1/2)
 -\sqrt x\,B_a\widetilde W(1/2)+O_{q,W}(1).
\]
The result holds also for the truncated rule in Corollary~\ref{cor:trunc} with $d=\varphi(q)$.
\end{corollary}
\begin{proof}
Character orthogonality expresses $S_a$ as the average of $\overline{\chi(a)}S_\chi$ for characters modulo $q$. Replacing an imprimitive character with its primitive inducing character changes $S_\chi(y)$ only at powers of the primes dividing $q$. For any $y>0$ there are at most two powers of a fixed prime in $[y,2y]$, so the change is $O_{q,W}(1)$ uniformly in $y$. The weighted average of this error is still $O_{q,W}(1)$ because $M_\mu(-1/2)<\infty$; the truncated moments are bounded uniformly as well. Now apply Theorem~\ref{thm:prime} and its truncated version to each primitive factor. Conjugation pairs make $B_a$ real. The only pole main term comes from the character induced by conductor $1$.
\end{proof}

\section{Examples, sharpness, and limitations}
\begin{center}
\begin{tabular}{@{}lll@{}}
\toprule
Family & $d$ & optimal exponent $\alpha$\\\midrule
$\zeta$ alone & 1 & 2\\
$L(s,\chi_{-4})$ alone & 1 & 2\\
$\zeta$ and $L(s,\chi_{-4})$ & 2 & 1\\
One nonreal character and its conjugate & 2 & 1\\
All primitive inducing characters modulo $q$ & $\varphi(q)$ & $2/\varphi(q)$\\
\bottomrule
\end{tabular}
\end{center}
In the first example the classical Riemann kernel \eqref{eq:xi}, normalized by $\xi(1/2)$, already gives the desired density; compact smoothing is unnecessary. For the primitive odd character modulo $4$, $\chi_{-4}(2j+1)=(-1)^j$, the theta series on the right of zero is alternating with strictly decreasing positive term magnitudes, since $y\ee^{-cy^2}$ is decreasing for $y\ge1$ when $c\ge\pi/4$. Its root number is $+1$, so its kernel is even and positive. In particular its central value is positive. The convolution of these two normalized positive kernels provides a concrete $d=2$ example without central division or any further smoothing. The supplied numerical script compares the theta integrals to the completed functions, tests cancellation at selected zero ordinates, and illustrates the general compact-smoothing lemma on a signed test density. It is a floating-point illustration, not a premise of the theorem.

A few distinctions are important. The theorem uses multiplicity-preserving cancellation: if two functions share an ordinate, their multiplicities are added. A weaker requirement of one vanishing at each distinct ordinate could admit a faster-decaying filter, since it need not pay for every zero counted with multiplicity. The exponent $2/d$ is sharp for the stated stronger requirement. In particular the lower bound does not depend on the zeros being distinct. A compactly supported measure would satisfy \eqref{eq:moment} for every $\alpha$, so exact global cancellation is impossible with compact support, in accordance with the finite-height obstruction in the merged article. The construction does not give explicit uniform conductor bounds for the constants or an efficient algorithm for determining central orders. No numerical central-zero decisions are used in its proof. The endpoint moment constant in Corollary~\ref{cor:constant} is only a necessary bound; optimality of that constant, and of the approximate finite-range width, remain separate questions.

Finally, this method does not extend the proved input to nonabelian Artin, general Hecke, or automorphic GRH. It uses the particular theta kernels, known zero counting, and the right-half-plane nonvanishing theorem for primitive Dirichlet functions. The positivity lemma itself is unconditional, and could be applied to other analytic completions only after verifying its precise signed-kernel and Laplace-positivity hypotheses. The article's proof of Dirichlet GRH supplies the input, not a claim that any new formalization or verification was performed here.

\begin{thebibliography}{9}
\bibitem{vendor} ProveIt, \emph{Vendored grh-lean (readable edition)}, upstream commit \texttt{72fead102f663d5c7a407964a13dc0cf9d977207}. \url{https://github.com/VladimirReshetnikov/ProveIt/tree/main/Analysis/GeneralizedRiemannHypothesis/vendor/grh-lean}.
\bibitem{article} ProveIt Contributors, \emph{Beyond the Riemann hypothesis for Dirichlet L-functions: consequences, transfers and limits}, merged article, 26 November 2026, especially sections \texttt{src}, \texttt{rec}, and \texttt{zer}. \url{https://github.com/VladimirReshetnikov/ProveIt/tree/main/Analysis/GeneralizedRiemannHypothesis/Article}.
\bibitem{davenport} H. Davenport, \emph{Multiplicative Number Theory}, 3rd ed., revised by H. L. Montgomery, Springer, 2000.
\bibitem{titchmarsh} E. C. Titchmarsh, \emph{The Theory of the Riemann Zeta-function}, 2nd ed., revised by D. R. Heath-Brown, Oxford University Press, 1986.
\bibitem{polson} N. G. Polson, \emph{The Riemann $\Xi$-Function as a Characteristic Function}, preprint, August 2026. \url{https://ssrn.com/abstract=7277559}.
\end{thebibliography}
\end{document}
